% 数学ノート集・共通テンプレート v2（積分論2026年度第1部の版面を参考）
% A5判。章内で定義・定理・例・注意等の番号を共通化。
% コンパイル： platex -interaction=nonstopmode math_note_template_v2.tex (2回)
%              dvipdfmx math_note_template_v2.dvi
\documentclass[a5paper,10pt,onecolumn]{jsarticle}
\usepackage[dvipdfmx,top=22mm,bottom=20mm,inner=19mm,outer=18mm,headheight=13pt,headsep=6mm]{geometry}
\usepackage{amsmath,amssymb,amsthm,mathtools}
\usepackage{enumitem,fancyhdr,titlesec}
\usepackage[dvipdfmx,bookmarks=true,colorlinks=false]{hyperref}
\usepackage{pxjahyper}

% ----- 差し替える文書情報 -----
\newcommand{\NoteTitle}{三次・四次方程式の解の公式とLagrange分解式}
\newcommand{\NoteSeries}{数学ノート集}

% ----- 文字・組版 -----
% 参考PDFの質素な明朝＋Computer Modernの印象を引き継ぎながら，
% A5判で窮屈にならないよう本文の行間を少し確保する．
\linespread{1.13}
\setlength{\parindent}{1zw}
\setlength{\parskip}{0pt}
\setlength{\abovedisplayskip}{6pt plus 2pt minus 2pt}
\setlength{\belowdisplayskip}{6pt plus 2pt minus 2pt}
\setlength{\abovedisplayshortskip}{3pt plus 2pt minus 1pt}
\setlength{\belowdisplayshortskip}{4pt plus 2pt minus 1pt}
\setlist[enumerate]{topsep=3pt,itemsep=1pt,parsep=0pt,leftmargin=2.4zw}
\titleformat{\section}{\normalfont\normalsize\bfseries}{\thesection}{0.8zw}{}
\titlespacing*{\section}{0pt}{2.0ex plus .4ex}{.9ex}
\titleformat{\subsection}{\normalfont\normalsize\bfseries}{\thesubsection}{0.7zw}{}
\titlespacing*{\subsection}{0pt}{1.5ex plus .3ex}{.6ex}
\raggedbottom

% ----- 柱とページ番号：ページ上部、罫線なし -----
\fancypagestyle{mathnote}{%
  \fancyhf{}%
  \fancyhead[L]{\footnotesize \NoteSeries\ （\NoteTitle）}%
  \fancyhead[R]{\footnotesize \thepage}%
  \renewcommand{\headrulewidth}{0pt}%
  \renewcommand{\footrulewidth}{0pt}%
}
\pagestyle{mathnote}

% ----- 節内通し番号の定義・定理群 -----
\newtheoremstyle{mathnotetheorem}
 {4.5pt}{4.5pt}{\normalfont}{}%
 {\bfseries}{}{0.65zw}%
 {\thmname{#1}\thmnumber{ #2}\thmnote{（#3）}}
\theoremstyle{mathnotetheorem}
\newtheorem{theorem}{定理}[section]
\newtheorem{definition}[theorem]{定義}
\newtheorem{lemma}[theorem]{補題}
\newtheorem{proposition}[theorem]{命題}
\newtheorem{corollary}[theorem]{系}
\newtheorem{example}[theorem]{例}
\newtheorem{remark}[theorem]{注意}
\renewcommand{\proofname}{証明}
\renewcommand{\qedsymbol}{$\square$}
\makeatletter
\renewenvironment{proof}[1][\proofname]{%
  \par\pushQED{\qed}\normalfont\topsep4\p@\@plus2\p@\relax
  \trivlist\item[\hskip\labelsep\bfseries #1]\ignorespaces
}{\popQED\endtrivlist\@endpefalse}
\makeatother

% ----- 使うノートで適宜拡張する記号 -----
\newcommand{\R}{\mathbb{R}}
\newcommand{\norm}[1]{\lVert #1\rVert}
\newcommand{\abs}[1]{\lvert #1\rvert}

\begin{document}
\section*{読了後にできること}
本稿を読むための前提は，三次方程式，判別式と平方根の選択．
本稿の具体的な到達課題は，三次方程式を抑圧三次式へ変換して根を構成する．

\section{方程式と根の対称式}

本稿では標数0の体 $K$ を考え，必要に応じて
その代数閉包 $\overline K$ を固定する．
解の公式は体拡大の中で根号を選択するという意味で扱う．

\begin{lemma}[Vietaの公式]\label{lem:vieta}
$K$ 上のモニック多項式が
\[
 f(X)=\prod_{j=1}^n(X-\alpha_j)
 =X^n+a_1X^{n-1}+\cdots+a_n
\]
と分解するとき，
\[
 a_k=(-1)^k\sum_{1\le i_1<\cdots<i_k\le n}
 \alpha_{i_1}\cdots\alpha_{i_k}.
\]
\end{lemma}
\begin{proof}
積 $\prod_j(X-\alpha_j)$ の $X^{n-k}$ の係数は，
$n$ 個の因子のうちちょうど $k$ 個から
$-\alpha_j$，残りから $X$ を選んだ積の総和である．
この選択は $k$ 元部分集合に一対一に対応するので
表示式が従う．
\end{proof}

\begin{proposition}[平行移動による二次項の消去]
\label{prop:depressed}
三次方程式
$X^3+aX^2+bX+c=0$ で
$X=Y-a/3$ と置くと
\[
 Y^3+pY+q=0,
 \quad p=b-\frac{a^2}3,\quad
 q=c-\frac{ab}3+\frac{2a^3}{27}.
\]
四次方程式
$X^4+aX^3+bX^2+cX+d=0$ は
$X=Y-a/4$ により
$Y^4+pY^2+qY+r=0$ へ帰着する．
\end{proposition}
\begin{proof}
二項定理を用いて $X=Y-a/3$ を三次式へ代入し，
$Y^2$ の係数が $-a+a=0$ となることを確認する．
残りの一次係数と定数係数を整理すれば
$p,q$ の式が得られる．
四次では $Y^3$ の係数は
$4(-a/4)+a=0$ であり，残る三係数を
$p,q,r$ と置けばよい．
\end{proof}

\section{Cardanoの公式}

\begin{theorem}[三次方程式の根号表示]\label{thm:cardano}
標数0の体 $K$ 上の三次方程式
\[
 Y^3+pY+q=0
\]
を考える．
\[
 \Delta=\left(\frac q2\right)^2+
          \left(\frac p3\right)^3
\]
と置き，$u^3=-q/2+\sqrt\Delta$，
$v^3=-q/2-\sqrt\Delta$，
$uv=-p/3$ を満たす $u,v\in\overline K$ を取ると，
$u+v$ は方程式の根である．
$\zeta^3=1$ なる原始3乗根 $\zeta$ を取れば，
三根は重複を込めて
\[
 u+v,\quad\zeta u+\zeta^2v,\quad
 \zeta^2u+\zeta v
\]
で与えられる．
\end{theorem}
\begin{proof}
$Y=u+v$ と置くと
\[
 Y^3+pY+q
 =u^3+v^3+(3uv+p)(u+v)+q.
\]
仮定から $u^3+v^3=-q$，
$3uv+p=0$ なので右辺は零である．
さらに
$\zeta^j u$ と $\zeta^{-j}v$ は
積 $uv$ とそれぞれの三乗を変えないので，
$j=0,1,2$ について全て根となる．
三次多項式の根は重複を込めてちょうど三個なので
これらが全ての根である．

$u,v$ の選択可能性を確認する．
$p\ne0$ なら
$U=-q/2+\sqrt\Delta$ を取り，
その三乗根 $u\ne0$ を選ぶ．
$v=-p/(3u)$ とすれば
\[
 v^3=-\frac{p^3}{27U}.
\]
$U(-q/2-\sqrt\Delta)
 =q^2/4-\Delta=-p^3/27$
なので $v^3=-q/2-\sqrt\Delta$ が従う．
$p=0$ なら $v=0$，$u^3=-q$ と選べばよい．
このとき上の式で $u+v=\zeta^ju$ の三根を得る．
\end{proof}

\begin{example}[二次根号の見掛け上の複素化]
$Y^3-3Y-1=0$ のCardano判別式は
$\Delta=1/4-1=-3/4<0$ である．
従って三実根を持つ場合でも，
公式の途中に複素平方根が現れ得る．
\end{example}
\begin{proof}
微分すると
$f'(Y)=3(Y^2-1)$ であり，
$f(-1)=1$，$f(1)=-3$ となる．
$f$ は負の無限遠で負，正の無限遠で正なので，
中間値の定理から
$(-\infty,-1)$，$(-1,1)$，$(1,\infty)$ に
それぞれ一つ根を持つ．
三次式なので三根は相異なる実数である．
一方 $\Delta=-3/4$ なので平方根は非実数となる．
\end{proof}

\section{三次のLagrange分解式}

\begin{definition}
三次方程式の根を $\alpha,\beta,\gamma$ とし，
原始3乗根 $\zeta$ を固定する．
\[
 L_1=\alpha+\zeta\beta+\zeta^2\gamma,\qquad
 L_2=\alpha+\zeta^2\beta+\zeta\gamma
\]
をLagrangeの分解式という．
\end{definition}

\begin{lemma}[分解式の三乗]\label{lem:lagrange-cubes}
三次式を $Y^3+pY+q$ とすると
\[
 L_1L_2=-3p,\qquad
 L_1^3+L_2^3=-27q,\qquad
 L_1^3L_2^3=-27p^3.
\]
\end{lemma}
\begin{proof}
根の和は $\alpha+\beta+\gamma=0$，
二根の積の和は
$\alpha\beta+\beta\gamma+\gamma\alpha=p$
である．
$\zeta+\zeta^2=-1$，
$\zeta^3=1$ を用いて
\begin{align*}
 L_1L_2
 &=\alpha^2+\beta^2+\gamma^2
   -(\alpha\beta+\beta\gamma+\gamma\alpha)\\
 &=(-2p)-p=-3p.
\end{align*}
また $L_1+L_2=3\alpha$ なので
\[
 L_1^3+L_2^3
 =(L_1+L_2)^3-3L_1L_2(L_1+L_2)
 =27\alpha^3+27p\alpha=-27q.
\]
最後の式は $L_1L_2=-3p$ を三乗して得る．
\end{proof}

\begin{proposition}[分解式と置換群]\label{prop:lagrange-action}
根の巡回置換
$(\alpha,\beta,\gamma)\mapsto
(\beta,\gamma,\alpha)$ は
\[
 (L_1,L_2)\longmapsto(\zeta^2L_1,\zeta L_2)
\]
を引き起こす．
二つの根の交換は $L_1,L_2$ を交換する．
従って $L_1^3,L_2^3$ は $A_3$ の作用で不変であり，
$S_3$ はその二つを置換する．
\end{proposition}
\begin{proof}
$\beta+\zeta\gamma+\zeta^2\alpha
=\zeta^2(\alpha+\zeta\beta+\zeta^2\gamma)$
であり，
$\beta+\zeta^2\gamma+\zeta\alpha=\zeta L_2$
である．
$\beta,\gamma$ を交換すると
$L_1=\alpha+\zeta\beta+\zeta^2\gamma$
は $L_2$ に変わり，逆も同様である．
三乗は $\zeta^3=1$ から巡回置換で不変であり，
$A_3$ は3-cycleで生成されるので結論を得る．
\end{proof}

\section{Ferrariの四次解法}

\begin{theorem}[四次方程式の二次式への分解]
\label{thm:ferrari}
四次式
\[
 f(Y)=Y^4+pY^2+qY+r
\]
に対して
\[
 R(Z)=Z^3+2pZ^2+(p^2-4r)Z-q^2
\]
と置く．$q\ne0$ とする．
$R(z)=0$ となる任意の根 $z$ は非零である．
$u^2=z$ を満たす $u$ を選び，
\[
 v=\frac{p+z-q/u}{2},\qquad
 w=\frac{p+z+q/u}{2}
\]
と置くと
\[
 f(Y)=(Y^2+uY+v)(Y^2-uY+w)
\]
となる．
従って四次方程式は三次方程式一つと
二次方程式二つを解くことに帰着する．
\end{theorem}
\begin{proof}
$R(0)=-q^2\ne0$ なので $z=0$ は根でない．
従って $u\ne0$ である．
右辺の二次式の積を展開すると
\[
 Y^4+(v+w-u^2)Y^2+u(w-v)Y+vw
\]
となる．
$v+w=p+u^2$，
$u(w-v)=q$ が定義から成り立つ．
また
\[
 4vw=(p+z)^2-\frac{q^2}z
\]
であり，
$R(z)=z(p+z)^2-4rz-q^2=0$ だから
$vw=r$ である．
従って係数は $p,q,r$ と一致して分解式が得られる．
三次式 $R$ はCardanoの公式で根号可解であり，
$u$ と二つの二次式の根は平方根で表される．
$q=0$ の場合は $Z=Y^2$ と置く二次方程式で
直接解ける．
\end{proof}

\section{四次方程式と三つの分割}

\begin{definition}
四次方程式の根を
$\alpha_1,\alpha_2,\alpha_3,\alpha_4$ とする．
四つの添字の二組への分割
\[
 \{12|34,\ 13|24,\ 14|23\}
\]
に対して
\[
 z_1=(\alpha_1+\alpha_2)^2,\quad
 z_2=(\alpha_1+\alpha_3)^2,\quad
 z_3=(\alpha_1+\alpha_4)^2
\]
と置く．
ここでは根の総和は零となるように
四次式の三次項を消去してあるとする．
\end{definition}

\begin{theorem}[三次分解式の根]\label{thm:resolvent}
上の $z_1,z_2,z_3$ は，
定理\ref{thm:ferrari}で定義した
\[
 R(Z)=Z^3+2pZ^2+(p^2-4r)Z-q^2
\]
の三根である．
\end{theorem}
\begin{proof}
$\alpha_1+\alpha_2=-(\alpha_3+\alpha_4)$
だから，四つの根を
\[
 \alpha_1,\alpha_2
 =\frac{-u\pm\sqrt{u^2-4v}}2,\qquad
 \alpha_3,\alpha_4
 =\frac{u\pm\sqrt{u^2-4w}}2
\]
と二組にまとめることができる．
ここで $u=-(\alpha_1+\alpha_2)$，
$v=\alpha_1\alpha_2$，
$w=\alpha_3\alpha_4$ と置いた．
Vietaの公式から
$v+w-u^2=p$，
$u(w-v)=q$，
$vw=r$ が成り立つ．
従って $z=u^2=(\alpha_1+\alpha_2)^2$
は $R(z)=0$ を満たす．
同じ議論は他の二組への分割にも適用できるので，
$z_1,z_2,z_3$ は全て $R$ の根である．
三つの値が重複しても，根を重複度を込めて考えれば
次数3の多項式の三根を与える．
\end{proof}

\begin{proposition}[四次の置換表現]\label{prop:s4-s3}
$S_4$ の三分割集合への作用は全射準同型
\[
 S_4\longrightarrow S_3
\]
を与え，その核は
\[
 V_4=\{1,(12)(34),(13)(24),(14)(23)\}
\]
に等しい．
\end{proposition}
\begin{proof}
$S_4$ は四元集合の二組への分割を置換するので
三元集合への置換表現が得られる．
$(12)$ は $13|24$ と $14|23$ を交換し，
$(23)$ は $12|34$ と $13|24$ を交換するため，
像は $S_3$ の互換を生成し全射である．
$S_4$ の位数は24，$S_3$ の位数は6だから
核の位数は4である．
$1$ と三つの二重互換は三分割をそれぞれ保存するので
核に含まれる．これらは四元あるから核に等しい．
\end{proof}

\section{判別式と置換の符号}

\begin{definition}
根 $\alpha_1,\dots,\alpha_n$ を持つ
モニック多項式の判別式を
\[
 \operatorname{Disc}(f)
 =\prod_{1\le i<j\le n}(\alpha_i-\alpha_j)^2
\]
と定める．
\end{definition}

\begin{proposition}[判別式の平方根]\label{prop:disc}
$f\in K[X]$ を標数0の分離的多項式とし，
分解体を $L$ とする．
\[
 \delta=\prod_{i<j}(\alpha_i-\alpha_j)
\]
と置くと，
任意の $\sigma\in\operatorname{Gal}(L/K)\le S_n$ に対して
\[
 \sigma(\delta)=\operatorname{sgn}(\sigma)\delta.
\]
従って
$\operatorname{Gal}(L/K)\subset A_n$ と
$\operatorname{Disc}(f)$ が $K$ の平方元であることは同値である．
\end{proposition}
\begin{proof}
置換 $\sigma$ は差の積を並べ替える．
一つの互換を適用すると積の符号が反転するので，
互換の積による表示から
$\sigma(\delta)=\operatorname{sgn}(\sigma)\delta$
が従う．
判別式は $\delta^2$ なので $K$ に属する．
$f$ は分離的だから $\delta\ne0$ である．
もしGalois群が $A_n$ に含まれれば
全ての自己同型が $\delta$ を固定するから，
Galois対応により $\delta\in K$ となり
判別式は平方元である．
逆に判別式が $b^2$（$b\in K$）なら，
体 $L$ で $(\delta-b)(\delta+b)=0$ だから
$\delta=\pm b$ となる．
従って $\delta\in K$ であり，
全ての $\sigma$ が $\delta$ を固定するため
符号は全て正である．
\end{proof}

\begin{remark}
三次・四次の解法は，根の置換作用に関して
適切に対称化した量を先に求めるという共通の構造を持つ．
三次のLagrange分解式は3-cycleで三乗が不変になり，
四次の三次分解式は $S_4/V_4\cong S_3$ の置換作用を
反映する．この群論的構造が，
三次・四次方程式の根号表示の根拠となる．
\end{remark}

\section{到達課題の完全解答}
\noindent\textbf{Cardanoの方法を実際に使う．}
\begin{proposition}
三次方程式 $x^3-3x-2=0$ の根は
$2,-1,-1$ である．
\end{proposition}
\begin{proof}
$x=2$ を代入すると $8-6-2=0$ であるから
$x-2$ で割り，
\[
x^3-3x-2=(x-2)(x^2+2x+1)
=(x-2)(x+1)^2
\]
を得る．
\end{proof}
\begin{proposition}[Cardanoの置換]
$x^3+px+q=0$ に対し $x=u+v$ と置き，
$uv=-p/3$ を課せば
\[
u^3+v^3=-q,\qquad
u^3v^3=-p^3/27.
\]
よって $u^3,v^3$ は
$t^2+qt-p^3/27=0$ の根である．
\end{proposition}
\begin{proof}
$(u+v)^3+p(u+v)+q
=u^3+v^3+(3uv+p)(u+v)+q$ を展開する．
指定した条件を使って整理すると表示式となる．
\end{proof}
\noindent\textbf{到達確認．}
$x^3-3x+2=0$ の因数分解を求めよ．
\par\noindent\textbf{解答．}
$x=1$ が根であり，
$x^3-3x+2=(x-1)^2(x+2)$ となる．
\end{document}